\documentclass[10pt,a4paper]{amsart}
\usepackage[margin=19mm]{geometry}
\usepackage{amsmath,amssymb,mathtools}
\usepackage[scr=rsfs,cal=euler]{mathalfa}
\usepackage{xfrac}
\usepackage[unicode,hidelinks,bookmarks=false]{hyperref}
\newcommand{\ie}{i.e.\ }
\newcommand{\cf}{cf.\ }
\newcommand{\resp}{respectively\ }
\newcommand{\Subsection}{\subsection}
\providecommand{\nobreakdash}{\nobreak\hbox{-}}
\setlength{\emergencystretch}{2em}
\multlinegap0pt
\usepackage{bm}
\usepackage{bbm}
\usepackage{dsfont}
\usepackage{xspace}
\usepackage{cancel}
\usepackage{mathtools}
\usepackage[shortlabels]{enumitem}
\usepackage{amsthm}
\makeatletter
\@ifundefined{lem}{\newtheorem{lem}{Lemma}}{}
\makeatother
\title[Integral quadratic forms over a ring of $p$-adic integers]{Integral quadratic forms over a ring of $p$-adic integers}
\author{Mrunal Hardikar, Anuradha S. Garge}
\date{Source: arXiv:2608.24808v1 (CC BY 4.0)}
\begin{document}
\maketitle

\noindent\emph{Source and licence.} Integral quadratic forms over a ring of $p$-adic integers. Version arXiv:2608.24808v1 (CC BY 4.0), distributed under the Creative Commons Attribution 4.0 International licence, as recorded in the arXiv licence metadata for this exact version and shown on the abstract page https://arxiv.org/abs/2608.24808v1. Full rights notice: licenses/arxiv-2608.24808-CC-BY-4.0.txt.\par\smallskip

\noindent\textit{Editorial scope.} Counted unit: the Lem whose source label is \texttt{Lemma 2.6 for odd prime} in the article (source0.tex lines 421--426, source label \texttt{Lemma 2.6 for odd prime}), delivered together with the article's own complete original proof of it (source0.tex lines 427--510, one \texttt{proof} environment of 84 lines). No mathematics is written, repaired or re-derived: statement and proof are the article's own text line for line. Required context retained uncounted: none. Every definition, display and label of those lines is retained unchanged. Omitted from this package and not redistributed: 1--420, 511--640. Nothing inside a retained excerpt is omitted or altered: every retained body is delivered exactly as the article's lines are written, so no omission receipt of any kind accompanies this package. Citations: the retained excerpts carry no citation command at all, so no bibliography entry of the article is transcribed and no reference is redistributed. Reference adaptation: none was needed and none was made; every reference inside the delivered unit resolves inside it, so no unresolved reference or citation is produced. Printed slips kept literal: no printed word, symbol, hypothesis or number of the counted statement or of its proof was altered. Layout and class: the article's own class is \texttt{article} with packages including fontenc, amsthm, amssymb, amsmath, stmaryrd, float, enumerate, array, setspace, cite, varioref, hyperref; the wrapper is the lane's pinned \texttt{amsart} editorial wrapper at 10pt on A4, which restates the theorem-like environments the retained text needs and adds the article's own macro definitions that the retained text needs (none), with extra wrapper packages (bm, bbm, dsfont, xspace, cancel, mathtools, enumitem). The delivered numbering is the unit's own. Assets: no retained span contains a figure environment, a graphics-inclusion command, a PDF-page inclusion, a table or a TikZ picture; no image file is copied into this package, the package declares no asset dependency and no image file is redistributed with it. Licence: version arXiv:2608.24808v1 of this article is distributed under the Creative Commons Attribution 4.0 International licence, as recorded in the arXiv licence metadata for this exact version and shown on the abstract page https://arxiv.org/abs/2608.24808v1. A full rights notice is delivered at licenses/arxiv-2608.24808-CC-BY-4.0.txt. Characters: every retained excerpt is pure ASCII, so the delivered engine, XeLaTeX with its default Latin Modern text font, reports no missing character.
	\begin{lem} \label{Lemma 2.6 for odd prime}
		Let $a_2$ and $a_3$ be units in $\mathbb{Z}_p$, $P \in M_{(n-2) \times 2}(\mathbb{Z}_p)$ and $Q \in M_2(\mathbb{Z}_p)$. Then there are $X_2$ and $X_3 \in M_n(\mathbb{Z}_p)$ such that $J^2 + \begin{bmatrix}
			0 & P\\
			0 & Q
		\end{bmatrix} - a_2X_2^2 = a_3 X_3^2 $. 
	\end{lem}

	\begin{proof}
		Since $a_2$ is unit in $\mathbb{Z}_p$, let $t_2 \in \mathbb{Z}_p$ such that $a_2t_2=1$. Let
		\begin{equation*}
			B_2 = s_1 E_{n-2,n} + s_2 E_{n-1,n} +s_3 E_{n,n}.
		\end{equation*}
		Let	$X_2 = J_{t_2} + B_2$. Hence
		\begin{eqnarray*}
		    X_2^2 &=& (J_{t_2} + B_2)^2 \\
			&=& J^2_{t_2} + J_{t_2} B_2 + B_2J_{t_2} +B_2^2 \\
			&=& t_2 J^2 + (J_{t_2} B_2 + B_2J_{t_2} +B_2^2)
		\end{eqnarray*}
		Now \begin{eqnarray*}
			J_2B_2 &=& s_1 E_{n-1,n} + t_2s_2 E_{n,n} \\
			B_2J_2 &=& (s_1 E_{n-2,n} + s_2E_{n-1,n} +s_3 E_{n,n}) (t_2 E_{n,n-1}) \\
			B_2^2 &=& s_1s_3 E_{n-2,n} + s_2s_3 E_{n-1,n} + s_3^2 E_{n,n}
		\end{eqnarray*}
		Since $a_2t_2=1$,
		\begin{equation} \label{Eq a2X2^2 for odd prime}
			a_2X_2^2 = J^2 +  \left[
			\begin{array}{c|cc}
				\textbf{0} & 0 & 0 \\ 
				\textbf{0} & s_1 & a_2s_1s_3 \\\hline
				\textbf{0} & s_2 & a_2 (s_1+s_2s_3) \\
				\textbf{0} & s_3 & s_2+ a_2s_3^2
			\end{array}
			\right]
		\end{equation}
		Let $L \in M_{(n-2) \times 2} (\mathbb{Z}_p)$ and $M = \begin{bmatrix}
			x & 1+x+x^2 \\
			-1 & -1-x
		\end{bmatrix}$
		Let $X_3 = \begin{bmatrix}
			0 & L \\
			0 & M
		\end{bmatrix}^2$. Note that $M^2 = \begin{bmatrix}
		-1-x & -1-x-x^2 \\
		1 & x
		\end{bmatrix}$ and $M^3 = I_2$. Therefore, $X_3 ^2 = \begin{bmatrix}
		0 & L \\
		0 & M
		\end{bmatrix}$. Hence
		\begin{equation} \label{Eq a3X3^2 for odd prime}
			a_3X_3^2 = \left[
			\begin{array}{c|cc}
				
				\textbf{0} & * &* \\ \hline
				\textbf{0} & a_3L_{n-2,1} & a_3L_{n-2,2} \\
				\textbf{0} & a_3x & a_3(1+x+x^2) \\
				\textbf{0} & -a_3 & a_3(-1-x)
			\end{array}
			\right]
		\end{equation}
		Let $Q = \begin{bmatrix}
			a & b \\
			c & d
		\end{bmatrix}$.
		From Equations \eqref{Eq a2X2^2 for odd prime} and \eqref{Eq a3X3^2 for odd prime},
		\begin{equation}
			J^2 + \begin{bmatrix}
				0 & P \\
				0 & Q
			\end{bmatrix} - a_2X_2^2 = a_3X_3^2
		\end{equation}
		if and only if the following equations are satisfied.
		\begin{subequations}
			\begin{align}
				a-s_2 &= a_3x  \label{Eq 10a} \\
				b- (a_2s_1 + a_2s_2s_3) &= a_3 (1+x+x^2) \label{Eq 10b}\\
				c-s_3 &= -a_3 \label{Eq 10c}\\
				d- (s_2 + a_2s_3^2) &= -a_3-a_3x \label{Eq 10d}\\
				P - \begin{bmatrix}
					0 & 0 \\
					s_1 & a_2s_1s_3
				\end{bmatrix} &= a_3L \label{Eq 10e}
			\end{align}
		\end{subequations}
		We can solve these equations for $x, s_1, s_2, s_3$ and $L$ as follows:
		\begin{enumerate}[(a)]
			\item From \eqref{Eq 10c}, find $s_3$.
			\item From \eqref{Eq 10a} and \eqref{Eq 10d}, $d-(s_2 + a_2s_3^2) = -a_3 -(a-s_2)$. Since $2$ is unit in $\mathbb{Z}_p$, we can find $s_2$. Since $a_3$ is unit, we can find $x$.
			\item Since $a_2$ is unit, we can find $s_1$ satisfying \eqref{Eq 10b}.
			\item Since $a_3$ is unit, we can find $L$ satisfying \eqref{Eq 10e}.
		\end{enumerate}
	\end{proof}

\end{document}
